\documentclass[11pt]{article}

% Change "review" to "final" to generate the final (sometimes called camera-ready) version.
% Change to "preprint" to generate a non-anonymous version with page numbers.
\usepackage[review]{acl}

% Standard package includes
\usepackage{times}
\usepackage{latexsym}

% For proper rendering and hyphenation of words containing Latin characters (including in bib files)
\usepackage[T1]{fontenc}
% For Vietnamese characters
% \usepackage[T5]{fontenc}
% See https://www.latex-project.org/help/documentation/encguide.pdf for other character sets

% This assumes your files are encoded as UTF8
\usepackage[utf8]{inputenc}

% This is not strictly necessary, and may be commented out,
% but it will improve the layout of the manuscript,
% and will typically save some space.
\usepackage{microtype}

% This is also not strictly necessary, and may be commented out.
% However, it will improve the aesthetics of text in
% the typewriter font.
\usepackage{inconsolata}

%Including images in your LaTeX document requires adding
%additional package(s)
\usepackage{graphicx}

\usepackage{enumitem}

%\usepackage{xurl}

% If the title and author information does not fit in the area allocated, uncomment the following
%
%\setlength\titlebox{<dim>}
%
% and set <dim> to something 5cm or larger.

\title{A Construct Validity Crisis in AI Sycophancy Evaluation}

% Author information can be set in various styles:
% For several authors from the same institution:
% \author{Author 1 \and ... \and Author n \\
%         Address line \\ ... \\ Address line}
% if the names do not fit well on one line use
%         Author 1 \\ {\bf Author 2} \\ ... \\ {\bf Author n} \\
% For authors from different institutions:
% \author{Author 1 \\ Address line \\  ... \\ Address line
%         \And  ... \And
%         Author n \\ Address line \\ ... \\ Address line}
% To start a separate ``row'' of authors use \AND, as in
% \author{Author 1 \\ Address line \\  ... \\ Address line
%         \AND
%         Author 2 \\ Address line \\ ... \\ Address line \And
%         Author 3 \\ Address line \\ ... \\ Address line}

\author{First Author \\
  Affiliation / Address line 1 \\
  Affiliation / Address line 2 \\
  Affiliation / Address line 3 \\
  \texttt{email@domain} \\\And
  Second Author \\
  Affiliation / Address line 1 \\
  Affiliation / Address line 2 \\
  Affiliation / Address line 3 \\
  \texttt{email@domain} \\}

%\author{
%  \textbf{First Author\textsuperscript{1}},
%  \textbf{Second Author\textsuperscript{1,2}},
%  \textbf{Third T. Author\textsuperscript{1}},
%  \textbf{Fourth Author\textsuperscript{1}},
%\\
%  \textbf{Fifth Author\textsuperscript{1,2}},
%  \textbf{Sixth Author\textsuperscript{1}},
%  \textbf{Seventh Author\textsuperscript{1}},
%  \textbf{Eighth Author \textsuperscript{1,2,3,4}},
%\\
%  \textbf{Ninth Author\textsuperscript{1}},
%  \textbf{Tenth Author\textsuperscript{1}},
%  \textbf{Eleventh E. Author\textsuperscript{1,2,3,4,5}},
%  \textbf{Twelfth Author\textsuperscript{1}},
%\\
%  \textbf{Thirteenth Author\textsuperscript{3}},
%  \textbf{Fourteenth F. Author\textsuperscript{2,4}},
%  \textbf{Fifteenth Author\textsuperscript{1}},
%  \textbf{Sixteenth Author\textsuperscript{1}},
%\\
%  \textbf{Seventeenth S. Author\textsuperscript{4,5}},
%  \textbf{Eighteenth Author\textsuperscript{3,4}},
%  \textbf{Nineteenth N. Author\textsuperscript{2,5}},
%  \textbf{Twentieth Author\textsuperscript{1}}
%\\
%\\
%  \textsuperscript{1}Affiliation 1,
%  \textsuperscript{2}Affiliation 2,
%  \textsuperscript{3}Affiliation 3,
%  \textsuperscript{4}Affiliation 4,
%  \textsuperscript{5}Affiliation 5
%\\
%  \small{
%    \textbf{Correspondence:} \href{mailto:email@domain}{email@domain}
%  }
%}

\begin{document}
\maketitle
\begin{abstract}
Sycophantic behavior is an established problem in AI systems, with concerns growing from users and researchers alike about the harmful effects of overly people-pleasing behavior in these models. However, there is disagreement between many researchers regarding what exactly counts as ``AI sycophancy'', and empirical studies primarily focus on a small subset of sycophantic behaviors with easily-measurable impacts on response quality. Further, few studies have evaluated the downstream impacts of sycophantic AI systems, particularly in the medium-to-long-term. In this position paper, we propose a unified definition of AI sycophancy, which involves displacement of context-relevant norms associated with approval-seeking behavior. We provide a framework for isolating, evaluating, and documenting sycophantic behaviors which prioritizes harm reduction and allows for a greater understanding of broader impacts. We outline a roadmap for this impact-centered framework, identify critical gaps in current evaluation practices, and provide recommendations regarding best practices and areas of future research. We believe that this framework is a step towards developing consensus, studying more diverse behaviors, and designing more effective evaluations for AI sycophancy.
% Think about EMNLP's ``Rethinking progress and evaluation in NLP'' goal - ``For example, going beyond static leaderboards, including real-world impact, trustworthiness, robustness, and longitudinal behavior; what “generalization” should mean for humans vs models, and when models are reused, adapted, and deployed in diverse settings, etc''
\end{abstract}

\section{Introduction}
There is growing concern regarding the phenomenon of ``AI sycophancy'' among researchers and users alike. A growing number of case studies have illustrated the potential catastrophic harms caused by excessive validation from AI systems [CITE], including hypotheses about potential links between sycophancy and an emerging phenomenon known as ``AI psychosis''. The critical nature of this topic necessitates careful empirical evaluations of these diverse phenomena and their downstream impacts. However, many AI sycophancy works are focused on a small subset of sycophantic behaviors constrained within a particular setting. There is a lack of data on the broader impacts of different forms of AI sycophancy, particularly the medium-to-long-term impacts. Further, there is a lack of consensus as to what ``counts'' as AI sycophancy among researchers, and no standardized definition of this term, which can result in broad generalizations being made about the nature of AI sycophancy and certain behaviors being overlooked in favor of those that are easier to evaluate. \citet{ye2026countsaisycophancytaxonomy} recently introduced a unified taxonomy for categorizing and discussing different types of AI sycophancy, which is a promising path towards standardizing the language we use to discuss AI sycophancy, and identified disagreements among researchers regarding the types of behaviors that should be considered sycophantic, referring to AI sycophancy as a \emph{fragmented construct}.

To address these issues in the field, we propose a unified, systematic, and normative approach to characterizing and evaluating AI sycophancy. We argue for a simple definition of AI sycophancy that aligns with the behaviorist nature of AI sycophancy research and promotes explorations of downstream harms of people-pleasing behaviors in AI systems. This definition can be formalized as \emph{user-approval-driven norm displacement}: a model behavior that accommodates a user’s belief, affect, identity, goal, or framing in order to preserve approval or relational ease, while reducing fidelity to a context-relevant norm such as truth, calibration, autonomy, care, fairness, task quality, or institutional integrity. We believe that this definition, and the accompanying framework we introduce, promote open investigation and documentation of the diverse behaviors in AI systems that optimize short-term user preferences over long-term outcomes. 

Our contributions can be listed as follows:

\begin{enumerate}
    \item Propose an empirical approach to characterizing AI sycophancy, which can resolve disagreements among researchers as to the types of behaviors that constitute sycophancy.
    \item Introduce a step-by-step framework for characterizing and evaluating AI sycophancy that centers impact and promotes the evaluation of a diverse set of behaviors
    \item Provide a roadmap for practitioners, which includes recommended best practices and gaps in the literature to inform future research
\end{enumerate}

% \begin{enumerate}
%     \item Motivation: findings related to people-pleasing in LLMs and case studies showing the harms of sycophancy
%     \item Trends in AI research, news coverage, and online discourse
%     \item Position: how researchers measure AI sycophancy often diverges from users' perceptions of AI sycophancy and major concerns about it.
%     \item Position: we need a more impact-centered framework for characterizing and evaluating AI sycophancy
%     \begin{enumerate}
%         \item Can help us better understand the dangers of this phenomenon, and which sycophantic behaviors causes the most harm, so we can target our mitigation strategies
%         \item Can also help us distinguish sycophancy from things like personalization and empathetic behavior - the line between adjusting for human preferences and excessive people-pleasing could be drawn based on the negative impacts over time
%     \end{enumerate}
%     \item One way to distinguish sycophancy from normal preference optimization, in the case of a lack of ground truth, could be by studying the impacts of this behavior (characterize ``excessive'')
%     \item Roadmap for the rest of the paper: discuss existing literature, argue for the need for more information regarding impacts of AI sycophancy, and enumerate recommendations for future work in this space
%     \end{enumerate}

    
\section{What exactly is 
``AI sycophancy''?}
The term ``sycophancy'' has long-standing origins, and is thought to originate from the Ancient Greek term, ``tale-teller about figs'' \cite{Theoddes64:online}. This term has come to refer to excessive people-pleasing behavior of various forms, often for personal gain; Merriam Webster defines sycophancy as ``obsequious flattery''\footnote{\url https://www.merriam-webster.com/dictionary/sycophancy} and a sycophant as ``one who praises those in power in order to gain their approval'', while the Cambridge Dictionary defines sycophancy as ``behavior in which someone praises powerful or rich people in a way that is not sincere, usually in order to get some advantage from them'' \footnote{\url https://dictionary.cambridge.org/us/dictionary/english/sycophancy}. 

In practice, ``sycophancy'' is typically used in a negative manner to characterize those who aim to deceive or ingratiate themselves for personal gain. This presupposes some notion of intent. On the contrary, AI sycophancy research primarily takes a behaviorist standpoint, quantifying this phenomenon as excessive people-pleasing in AI systems that negatively impacts their outputs. Different works conceptualize ``excessive'' in different ways, and many works focus on accuracy or factuality to quantify the negative impacts of people-pleasing behavior [CITE]. Without a notion of ground truth, drawing this distinction becomes less trivial. A few works utilize normative frameworks including Bayesian rationality \cite{atwell2025basil}, while other works characterize ``excessive'' people-pleasing behavior as deviations from existing evidence: for instance, overturning negative paper reviews in response to user rebuttals, even when unsubstantiated \cite{wang2026ai}. Still other works define sycophancy as people-pleasing behaviors that lead to inconsistencies in LLMs' stated positions or outputs. 

This divergence between definitions reflects a fundamental disagreement among researchers regarding the types of model behaviors that constitute sycophancy. \citet{ye2026countsaisycophancytaxonomy} recently conducted a large-scale literature review and surveyed over 100 experts on AI sycophancy, and observed substantial variation in how sycophancy is defined and measured by researchers. The authors differentiate between position sycophancy, where models may shift their stated positions to align with or flatter users, and person sycophancy, where models may praise or positively evaluate users' traits and validate their emotions. They find the latter to be underrepresented in AI sycophancy research, with a lack of consensus among researchers regarding whether implicitly expressed Person behaviors should be considered sycophantic. A few works detail specific instances in which the lines may be blurred between optimizing for user preferences and sycophancy in AI systems. In particular, the line between personalization and sycophancy has been debated in the context of political discourse \cite{batzner2025germanpartiesqa}, as has the distinction between empathy and sycophancy \cite{PDFAIEmp14:online}. 

\section{AI sycophancy as norm displacement} One thing that many of these divergent definitions have in common, beyond the characterization of sycophancy as people-pleasing behavior, is a negative view of these behaviors or downstream impacts. Sometimes this is implicit, with references to AI sycophancy as an alignment failure \cite{Davalas2026MitigatingSA} or use of the term ``excessive'' in their characterizations of AI sycophancy \cite{Lee2026EffectsOA}. Sometimes these definitions are attached to certain negative impacts; for instance, many papers specifically define AI sycophancy as the tendency to conform to user opinions at the cost of accuracy 
\cite{Ibrahim2026TrainingLM} or in contrast with existing evidence \cite{Yao2026HearingIB, wang2026ai, Shah2026SycoPhantasyQS}. Other negative impacts mentioned in existing definitions of sycophancy include ``reinforcing biases'' \cite{lahlou2026writing}, ``providing false objectivity to transient anxieties'' \cite{PDFAIEmp14:online}, and encouraging behavior or beliefs that are ``aligned with a user's short-term preferences'' but ``detrimental to their long-term interests'' \cite{fulay2026delegates}. Others define sycophancy more broadly as the tendency to prioritize user approval above ``truth'' \cite{sharma2024towards, turner2026programmed}, which may generalize beyond notions of ground truth; \citet{turner2026programmed}, for instance, characterizes ``truth'' as ``an orientation toward 'getting things right''', not limited to specific truth values. Some researchers implicitly distinguish sycophancy as a negative side-effect of optimizing for user preferences; in other words, for outputs to be sycophantic, there must be a negative impact on output quality or negative downstream impacts that could emerge as a result of demonstrated people pleasing behaviors. For instance, \citet{PDFAIEmp14:online} argue that what makes this behavior sycophantic is the negative impact on certain users, specifically younger users.

Traditional notions of sycophancy in human behavior are fundamentally incompatible with the behaviorist lens through which AI sycophancy is viewed, and these two notions cannot be reconciled without, to some degree, anthropomorphizing AI systems. As a result, AI sycophancy researchers fundamentally disagree on how this concept should be defined and operationalized \cite{ye2026countsaisycophancytaxonomy}, with different works characterizing this phenomenon in wildly different ways. We argue that these divergent definitions can be reconciled by characterizing AI sycophancy broadly as \textit{model behaviors that accommodate a user's perceived belief, affect, identity, goal, or framing, in order to preserve approval or relational ease, while reducing fidelity to a context relevant norm}. In other words, rather than use the displacement of specific norms to characterize AI sycophancy (factuality, ``truth'', fairness, etc.), we define it as a general norm violation caused by apparent approval-seeking behavior in AI systems. This allows us to empirically verify which behaviors should fall under the umbrella of AI sycophancy, based on whether they reduce fidelity to context-dependent norms. This definition also accommodates significant variation regarding the types of behaviors researchers consider ``sycophantic'', while addressing growing concerns among researchers regarding one-size-fits-all approaches for labeling and addressing AI sycophancy \cite{Vennemeyer2025SycophancyIN, ye2026countsaisycophancytaxonomy}.



% \begin{enumerate}
%     \item 
%     \item Sycophancy research in AI
%     \begin{enumerate}
%         \item Origin of the popularization of ``AI sycophancy'', and metrics showing the trend in the number of sycophancy mentions
%         \item How, and why, AI sycophancy diverges from the dictionary definition
%         \item Overview of AI sycophancy research and proposed taxonomy (cite Meryl's paper)
%         \item Studied domains
%         \item Studied modalities
%         \item Studied mitigation strategies
%         \item Evaluation metrics
%     \end{enumerate}
%     \item AI sycophancy research in other domains
%     \begin{enumerate}
%         \item Psychology
%         \item Philosophy
%     \end{enumerate}
% \end{enumerate}
\section{What do we know about the impacts of AI sycophancy?}
In order to use norm displacement as a dividing line between people-pleasing behaviors and sycophancy in LLMs, it is critical that researchers understand the normative effects of these behaviors: both the direct impacts on output quality and the downstream impacts in the short-, medium-, and long-term. We detail these impacts in the following paragraphs.

\subsection{Direct Impacts on Output Quality}
Most of the existing works on AI sycophancy have sought to quantify the direct impacts of apparent approval-seeking behavior on output quality. Some of these impacts are more trivial to measure, like changes in accuracy \cite{sharma2024towards}, calibration \cite{sicilia-etal-2025-accounting} and hallucinations \cite{sharma2024towards,rrv-etal-2024-chaos}. Other impacts may be more difficult to measure: for instance, contradicting existing evidence in scenarios that may be subjective or where the model has incomplete information \cite{atwell2025basil}, or encouraging clinically harmful patterns in mental health support scenarios \cite{han2026detecting}.

Some works also study these impacts in more specific contexts, domains, or applications.
\citet{wang2026ai} find that AI reviewers were vulnerable to reversing their positions when presented with user rebuttals, even if this reversal was unfounded and not supported by the given evidence. \citet{lahlou2026writing} found that AI-generated literature reviews contained biases reflecting the users' perceived viewpoints or intended directions. 


\subsection{Downstream Impacts}
Few existing papers have identified  downstream impacts of AI sycophancy, and those that do have primarily focused on short-term impacts. We describe these works below, distinguishing between two categories of sycophantic behaviors defined in \citet{ye2026countsaisycophancytaxonomy}: \emph{Position} and \emph{Person}, where Position sycophancy involves conforming to or validating a user's opinions or beliefs, and Person sycophancy involves validating the person themselves (including character traits and emotions). 

\paragraph{Person sycophancy}
While many works speculate on the impacts of Person sycophancy on humans' psychology and social well-being, only a few works quantify these directly. \cite{lee2026effects} study the impact of sycophancy on social well-being on the context of AI companionship, and find a negative relationship between sycophantic model responses and perceived social support. \cite{chandra2026sycophantic} model the connection between AI sycophancy and AI psychosis, and find that even ideal Bayesian reasoners can experience delusional spiraling as a result of chatbot conversations, and that sycophancy plays a role in this phenomenon. \cite{Batista2026ARA} study how sycophancy can distort users' reality and reinforce user biases, even when they contradict the existing evidence. \cite{rathje2025sycophantic} study the effects of sycophancy in the context of political conversations, and find evidence indicating that humans prefer sycophantic models, and that engaging with sycophantic models in these discussions increased attitude extremity and certainty, indicating the possibility that sycophantic AI systems may contribute to echo chamber formation. \citet{cheng2026sycophantic} find that sycophantic responses in advice-seeking conversations regarding interpersonal conflicts lead to users feeling less inclined to repair relationships and more certain that they are in the right. They also found that users prefer sycophantic responses overall, despite the evidence that they impaired the users' social judgments. 

All of these works provide evidence that Person sycophancy may have a profound impact on humans' psychological health and social well-being, but there are not yet any works that capture this impact beyond the short-term (within an isolated set of conversations). Given the mounting evidence of AI psychosis and psychological harm stemming from long-term interactions with AI chatbots \cite{TheyAske50:online}, it is crucial that researchers conduct more medium-to-long-term studies on the impact of AI sycophancy on users' mental wellbeing, particularly in naturalistic settings. There is also a critical need to study how AI sycophancy may impact developing brains, as features which may help adult users may harm youth. For instance, \citet{jiao2026ai} point out that emotional mirroring, a behavior rewarded by adult users, can promote emotional dependency in younger users and impair emotional development by reducing social friction necessary to facilitate emotional development and develop executive functions. 

\paragraph{Position sycophancy}
There is also a growing body of work demonstrating the results of Position sycophancy (where the user's position is validated) on task success in various domains.

A few works have studied the impact of Position sycophancy in decision-making settings. \cite{li2026does} found that opinion conformity reduced the likelihood of users changing their mind when opposing evidence was provided. \citet{sicilia-etal-2025-accounting} found that sycophancy can reduce the success of human-AI collaboration for uncalibrated users. \citet{kumar2025ai} found evidence that AI software engineering agents often did not push back on or debate human developers, and that when they were directly encouraged to do so, developers uncovered new insights about their codebases. 

A couple of works have also studied the impact of Position sycophancy in multi-agent systems. \citet{kasprova2026too} found that providing agents with sycophancy priors in collaborative multi-agent systems improved discussion accuracy by 10.5\%. \citet{Yao2025PeacemakerOT} find that, in the setting of multi-agent debate, sycophancy amplifies disagreement collapse before reaching a correct conclusion and results in lower accuracy than single-debate settings. 

These works represent instances in which researchers have isolated the negative impacts of Position sycophancy in task-specific settings. A much larger body of literature exists which identifies sycophancy in various domains (mental health, finance, healthcare and medicine, software, etc.) but significantly fewer works have studied its impact on human-AI collaboration and task success. Given the critical, high-stakes nature of many of these domains, such as law, healthcare, and education, there is a need to conduct more assessments regarding the impact of AI sycophancy on specific processes within these domains. Further, many of the above works have studied the impacts of sycophancy on the micro level, isolating specific instances and interactions in which sycophancy has an impact on task success. Upon identifying these impacts at a granular level, it is crucial that we understand the systematic and institutional effects of AI sycophancy: its impacts on the peer-review process, patient care, etc.

\section{The case for a unified construct}

Many researchers characterize AI sycophancy as an inherently negative property, despite the critical gaps in the field's understanding about the downstream effects of people-pleasing behavior in LLMs. This makes it difficult to determine which approval-seeking behaviors are associated with negative impacts, and should thus fall under the umbrella of ``AI sycophancy''. As a result, most researchers focus on studying opinion shifts, particularly for objective tasks, as accuracy changes, factual errors, and hallucinations are easy to identify. However, recent works have indicated that other, more under-studied sycophantic behaviors, including social sycophancy \cite{cheng2026sycophantic}, can have negative downstream effects on users.  %, but these behaviors are comparatively under-studied in current research.

Many works on AI sycophancy are also motivated by its potential long-term effects, but these effects are largely unstudied and unverified. Some critical failures of generative AI systems have been posited to result from sycophancy in AI, but it is typically unclear what, if any, sycophantic behaviors are connected to these impacts.
Without a clear understanding of the downstream effects of a diverse set of sycophantic or people-pleasing behaviors in AI systems, it is difficult for researchers to make informed decisions about which behaviors are most important to quantify and mitigate to prevent certain harms.

We propose a standardized, bottom-up approach to characterizing and evaluating sycophancy in AI systems, where the focus is on identifying and quantifying norm displacements. This aligns with the behaviorist framing of existing research in AI sycophancy and serves as a means to resolve disagreements among researchers about the types of behaviors that should be characterized as sycophantic. It also encourages more careful, data-driven benchmark development that measures a diversity of sycophantic behaviors associated with downstream harms. Further, it gives researchers a clearer intuition regarding the possible broader impacts of different sycophancy mitigation strategies, allowing them to deploy selected mitigation strategies to target specific harms.

This framework is also in alignment with broader calls in the NLP and machine learning communities to consider downstream impacts when designing model evaluations. \citet{paullada2021data} have called out a number of problems associated with dataset development and benchmarking in AI disciplines, including  \citet{eriksson2025can} recently pointed out a number of problems with current AI benchmarks, including misalignments between what benchmarks actually studied and what they were claimed to study, a lack of a clear definition of the phenomena that benchmarks were designed to measure, lack of considerations for downstream impacts, and a lack of clarity regarding how benchmark results should be used in practice. 

% \begin{enumerate}[noitemsep]
%     \item Focus research towards studying sycophantic behaviors that are associated with the most negative impacts
%     \item Develop benchmarks that are informed by a more complete understanding of the behaviors that are most harmful
%     \item Evaluate targeted mitigation strategies, and understand which mitigation strategies are best for addressing which downstream outcomes
% \end{enumerate}
\input{sections/04-impacts}
\begin{figure*}
    \centering
    \includegraphics[width=\linewidth]{figures/roadmap.png}
    \caption{A visualization of our proposed impact-centered framework for characterizing and evaluating AI sycophancy.}
    \label{fig:roadmap}
\end{figure*}

\section{A roadmap for impact-centered research}
In order to come to a consensus regarding which people-pleasing behaviors should be considered sycophantic in AI systems, what the far-reaching impacts of these behaviors are, and which mitigation strategies should be deployed to address these impacts, we propose a roadmap for identifying and investigating different forms of AI sycophancy and developing target benchmarks based on these findings. The steps outlined below do not need to be entirely linear; for instance, behaviors that have a measurable negative impact in the immediate term (e.g. on accuracy) can be identified as sycophantic even without fully understanding their downstream impacts. Similarly, researchers may succeed identifying certain behaviors as sycophancy based on their negative short-term impacts, but may want to investigate further regarding their medium-to-long term impacts. Different people-pleasing behaviors can also be studied in parallel, as opposed to making sweeping claims about AI sycophancy as a whole. We outline these steps below and, in the remaining paragraphs, provide suggestions, best practices, and future directions for each of these steps. These steps are illustrated in Figure \ref{fig:roadmap}.

\begin{enumerate}[noitemsep]
    \item Identify approval-seeking behaviors in LLMs (whether or not there is a consensus around these behaviors being sycophantic) 
    \item Study the impacts of these behaviors (direct impacts on output quality, as well as short-, medium-, and long-term downstream impacts) with respect to context-relevant norms. Report these impacts using a standardized framework (which we propose in \S\ref{sec:impact-card}) in order to better compare and aggregate the results from different studies.
    \item Isolate the people-pleasing behaviors associated with norm displacement in the immediate, short, medium, or long-term. These behaviors can be considered ``sycophantic''.
    \item Develop new benchmarks for these sycophantic behaviors, informed by knowledge of the characteristics that are particularly associated with negative impacts
    \item Evaluate different language models and mitigation strategies on this updated suite of benchmarks.
    \item Assess the short, medium, and long-term impacts of different baselines and mitigation strategies. Report the results of mitigation strategies using our Sycophancy Impact Cards proposed in \S\ref{sec:impact-card}.
\end{enumerate}

\noindent In following this roadmap, researchers should be considerate of the following:

\paragraph{Specificity} Researchers' notion of AI sycophancy comprises a broad set of behaviors, which may have different impacts on different domains or settings \cite{ye2026countsaisycophancytaxonomy}. These behaviors must be investigated individually in order to avoid adopting a one-size-fits-all approach to a diverse set of behaviors that fall under the umbrella of ``AI sycophancy''. \citet{ye2026countsaisycophancytaxonomy} develop a taxonomy which can serve as a starting point for researchers to categorize the specific behaviors they're testing.

\paragraph{Diversity} Many AI sycophancy studies focus on opinion conformity in LLMs [CITE], as this is one of the simplest approaches to directly measure and researchers can easily study degradations in quality in scenarios with a ground truth. However, many other sycophantic behaviors exist and may have the potential to harm users or negatively impact outcomes: for instance, excess flattery or uncritical emotional validation.

\subsection{Sycophancy Impact Card}
\label{sec:impact-card}
To develop a unified approach to quantifying and categorizing AI sycophancy along multiple dimensions, we propose that researchers create a \emph{Sycophancy Impact Card} to report the results of their works. This should include the following: 

\begin{enumerate}
    \item \textbf{Behavior(s) measured}: the types of approval-seeking behaviors observed in LLMs. Examples include \emph{opinion conformity} \cite{sharma2024towards, sicilia-etal-2025-accounting}, excessive \emph{flattery} \cite{sharma2024towards}, \emph{validation} \cite{cheng2025elephant}, and \emph{hedging} [CITE].
    \item \textbf{Referent}: From a comprehensive survey of AI sycophancy literature, \citet{ye2026countsaisycophancytaxonomy} find that documented sycophantic behaviors in AI are expressed towards either a user's expressed position(s) (\emph{Position}) or the user themselves (\emph{Person}), with most AI sycophancy works evaluating Position sycophancy. We recommend that authors adopt this binary label in order to categorize the behaviors studied in a standardized manner. The authors also distinguish sub-referents for each referent, to denote the explicitness of the sycophantic behavior. For Person sycophancy, these sub-referents are Traits vs. Emotions of a person, and for Position sycophancy, these sub-referents are Verifiable vs. Subjective position. We recommend that authors adopt this taxonomy as well.
    \item \textbf{Norm displaced}: Here, researchers should characterize the specific norm(s) that are found to be displaced in their particular study. Examples include \emph{accuracy} \cite{sharma2024towards}, factuality \cite{rrv-etal-2024-chaos}, rationality \cite{atwell2025basil, Batista2026ARA}
    \item \textbf{Outcome measured}: This should contain specifics about the outcome(s) that were measured for which norm displacement was observed. Examples include changes in multiple-choice question-answering performance on MMLU [CITE], willingness from users to repair relationships \cite{cheng2026sycophantic}, and user confidence in the outcome of a conversation \cite{sicilia-etal-2025-accounting}.
    \item \textbf{Time horizon}: the time period over which the outcome is being measured. The impacts of approval-seeking behavior on output quality would represent an immediate outcome, whereas the impacts over a single conversation (or multiple conversations within a short time span of minutes to hours) would represent short-term outcomes, the impacts over multiple conversations within a period of days to weeks would constitute medium-term outcomes, and the impacts over multiple conversations within a period of months to use would constitute long-term outcomes. 
    \item \textbf{Population} Studies should also document the particular population(s) on which they are studying, as sycophancy may have different outcomes on different demographics (for instance, young children vs. adults or more vs. less suggestible people) or different occupations (e.g. medicine vs. law).
    \item \textbf{Evidence level}: the \emph{evidence level} refers to the degree to which approval seeking behaviors (and their downstream outcomes) are being assessed. These levels are defined in our Evidence Maturity Ladder described in \S\ref{sec:evidence-maturity}. 
    \item \textbf{Proposed mitigations} If the paper evaluates the impacts of any proposed sycophancy mitigation techniques, these should be named and described here. 
    \item \textbf{Outcomes of proposed mitigations} If there are any proposed mitigations, the outcomes of each proposed mitigation technique should be detailed here.
\end{enumerate}

Below, we describe the different levels within the \emph{evidence maturity ladder}.

\subsection{Evidence Maturity Ladder}
\label{sec:evidence-maturity}

Our proposed evidence maturity ladder allows for easy categorization of the degree to which approval-seeking behaviors and their downstream outcomes are assessed. This ranges from behavioral detection (for instance, models switching to the user's answer from a different answer) to institutional and societal impacts. We describe each of these levels and provide some examples for each level below.

\begin{enumerate}
    \item \textbf{Behavioral detection}: identifying and studying an approval-seeking behavior in LLMs without quantifying any resulting norm displacements or downstream impacts falls under the first level of behavioral detection. 
    \item \textbf{Norm displacement}: works that fall under this level evaluate the extent to which specific context-dependent norms are displaced in the model's output. This could mean changes in accuracy, rationality, social biases, etc. in the generated response. The most common example in current research is measuring the effects of switching to the user's position on accuracy of model responses.
    \item \textbf{Interaction impact}: works that fall under this category study the impact of approval-seeking behaviors within the scope of a single interaction. Examples of this include \citet{cheng2026sycophantic} and \citet{rathje2025sycophantic}.
    \item {Repeated interaction impact}: works that study the impact of approval-seeking behaviors in AI over the course of multiple conversations fall under this level: for instance, asking users to take a psychological assessment before and after using a generative AI model for therapy for several weeks.
    \item \textbf{Institutional/societal impact} many works that discuss the problems with current benchmarks mention the lack of works that evaluate the larger-scale impacts of AI systems on institutions and society. Possible institutional impacts could include the effects of AI sycophancy on the overall peer review process, given that AI sycophancy has been shown to steer models towards over-correction in the face of rebuttals \cite{wang2026ai}. Another example is the impact on democracy, which \citet{turner2026programmed} argue may be negatively impacted by sycophantic behaviors.
\end{enumerate} 

\section{Towards a unified evaluation framework}

To assist future researchers with operationalizing our proposed Sycophancy Impact Card, we provide several case studies and examples of possible impact cards for existing works. We also release a GitHub repository with sycophancy impact cards for existing literature, which other researchers in the community can submit contributions to. This will be attached to a HuggingFace space with aggregated information regarding different sycophantic behaviors, documented norm displacements, and observed outcomes, in order to facilitate future exploration on AI sycophancy.

\subsection{Case Studies}

\paragraph{\citet{sharma2024towards}}: One of the foundational works in AI sycophancy, this work was largely responsible for coining the term \emph{AI sycophancy}, and the authors develop multiple novel benchmarks for studying this phenomenon. Multiple sycophantic behaviors are quantified here, including conformity to users' beliefs and subtly mimicking a user's mistake. We illustrate the corresponding impact card in Figure \ref{fig:impact-card}.

\begin{figure}
    \includegraphics[options]{figures/sharma2024_impact_card.pdf}
    \caption{Sycophancy impact card for \cite{sharma2024towards}}
    \label{fig:impact-card}
\end{figure}

\subsection{\citet{cheng2026sycophantic}}: This notable recent work that evaluates the short-term impacts of AI sycophancy studies the impacts of social sycophancy on human users after one advice-giving conversation. This work finds that users who chat with the sycophantic model are more likely to be convinced that they are right in interpersonal conflicts and less likely to repair relationships. Despite these negative effects, users are also shown to prefer the sycophantic responses overall. In Figure \ref{fig:impact-card-cheng-et-al}, we illustrate a sample impact card for this work. 

\begin{figure}
    \includegraphics[]{figures/cheng2026_impact_card.pdf}
    \caption{Sycophancy impact card for \cite{cheng2026sycophantic}}
    \label{fig:impact-card-cheng-et-al}
\end{figure}

% \subsection{Behavioral detection}

% Of the existing AI sycophancy literature, most empirical works fall, at least partially, under this step. People-pleasing behaviors that have been identified include opinion conformity [CITE], hedging [CITE], praise [CITE], emotional validation [CITE], and encouragement of users' actions [CITE]. People-pleasing can also manifest as a deficit of certain behaviors, for instance a lack of pushback [CITE]. 

% There are also some documented instances of people-pleasing that many researchers do not associate with sycophancy: for instance, personalization \cite{Batzner2024GermanPartiesQABC} and politeness \cite{turner2026programmed}. These are often subject to disagreement among researchers: for instance, \cite{10.1145/3772318.3790934} consider self-deprecation to be a dimension of AI sycophancy and include these behaviors in their analysis, while \citet{turner2026programmed} argue that this should be classified as epistemic humility rather than sycophancy. Steps 2 and 3 of our proposed roadmap can help researchers more clearly delineate these distinctions, and characterize the cases in which these behaviors should (or shouldn't) be considered ``sycophancy''. 

% \subsection{Norm displacement}

% Some works on AI sycophancy have been able to quantify the immediate impacts of certain people-pleasing behaviors on output quality, allowing them to more definitively support their claims that these behaviors should be considered sycophantic. The most striking example is opinion conformity, particularly in instances where there is a ground truth: many works have found a consistent pattern where LLMs are inclined to agree with users' judgments at the expense of accuracy [CITE]. Further, a few works have demonstrated the tendency for LLMs to hallucinate in order to align with users' perceived goals or beliefs \cite{rrv-etal-2024-chaos, sharma2024towards}. 

% Not all negative impacts on output quality can be assessed this way. Assessing negative impacts on users' psychological states and social interactions requires user testing and careful study design. Thus, only a few works have studied these impacts directly, and these works have primarily focused on short-term outcomes \cite{rathje2025sycophantic, lee2026effects, cheng2026sycophantic}. A few papers have also studied the impacts of AI sycophancy on certain processes or tasks (e.g. pair coding, multi-agent systems, peer review), but these have primarily been studied at an individual level; the effect of AI sycophancy on different institutions and domains at large is yet to be studied. Further, the medium-to-long-term impacts of AI sycophancy on task success, human decision-making, and user capabilities in cases of frequent, sustained human-AI interaction are yet to be explored. 

% Below, we outline some general recommendations for assessing the impacts of people-pleasing behaviors in AI:

% \paragraph{Data collection} Works that study the short-term impacts of AI sycophancy often do so at the conversation level, either with more controlled studies or naturalistic interactions \cite{rathje2025sycophantic, cheng2026sycophantic}. There is a need for works that study the impact of AI sycophancy over time and after sustained usage. This can be done by conducting studies in which users utilize AI systems over multiple sessions \cite{kosmyna2025your} or by observing naturalistic interactions with LLMs  over time. It is of note that these studies should sample a variety of users with different levels of familiarity or experience with AI, trust in AI systems, and vulnerability to deception, as these factors could moderate the impacts of sycophantic behaviors in different ways.
% \paragraph{Experimental settings} Prior work has tended towards either A/B testing (testing a "more sycophantic" model against a "less sycophantic" model) [CITE] or identifying specific sycophantic behaviors and measuring the association between these behaviors and the impacts of interest [CITE]. A/B testing allows researchers to make direct comparisons between the impacts of sycophantic vs. less-sycophantic models, whereas identifying specific sycophantic behaviors in human-AI interactions and measuring their association with individual conversational outcomes can allow researchers to build a causal model of the effects of certain people-pleasing behaviors [CITE].
% \paragraph{Evaluation} Evaluating the impacts on AI sycophancy will likely look different depending on the nature of the impacts of interest. Impacts on user psychology and social judgment should be directly assessed via questionnaires, and can also be indirectly assessed based on the dynamics of their conversations and frequency of interactions. Impacts on specific tasks or processed can be measured by quantifying overall task success or evaluating the quality of the final product(s) produced. Certain conversational dynamics can also be studied here: the degree to which humans and AI systems push back against one another, the frequency in which humans accept AI suggestions and vice versa, etc.

% \subsection{Isolate behaviors associated with negative impacts}
% What constitutes a negative impact?
% Accuracy, task success, etc. are easy to defend for machine learning researchers, but other impacts may be more subtle, such as an increase in cognitive offloading and a rise in homogeneity. AI sycophancy has been shown to promote dependence on language models \cite{cheng2026sycophantic}, which has been associated with reductions in linguistic diversity [CITE] and creative diversity [CITE]. It can be argued that these are considered a form of societal harm, as many existing works detail the importance of diversity of thought in collaborative settings, and users from marginalized groups whose writing styles are less represented in training data have been shown to be especially impacted [CITE]. 

% AI sycophancy researchers should also continue to study the psychological impacts of AI sycophancy. It has been shown that social sycophancy can reduce pro-social behavior in users \cite{cheng2026sycophantic}, but many other potential psychological impacts are underexplored. For instance, \citet{PDFAIEmp14:online} have raised the possibility that AI sycophancy may impact the development of key cognitive functions in younger users, including emotional management and critical thinking, by reducing the social friction necessary to develop these skills.

% It is possible that some people-pleasing behaviors in LLMs may simply be byproducts of preference modeling, without any negative or harmful consequences. These can be excluded from the umbrella of AI sycophancy, but should continue to be studied to rule out any downstream harms.

% \subsection{Develop new evaluations}
% Upon isolating new instances of sycophancy in LLMs and developing frameworks to assess their impacts, researchers should work backwards to develop new benchmarks that:
% \begin{enumerate}
%     \item \textbf{Benchmark a diverse set of people-pleasing behaviors that are proven to be linked negative downstream impacts}, like social sycophancy \cite{cheng2025elephant,cheng2026sycophantic}. Evidence of the negative impacts of a diverse set of people-pleasing behaviors should encourage more variety in benchmark development, rather than continuing to primarily focus on opinion conformity in scenarios with a ground truth, with accuracy as a metric.
%     \item \textbf{Assess the downstream impacts of different sycophantic behaviors in various domains, scenarios, cultural contexts, and applications spaces.} Once certain sycophantic behaviors are shown to have a negative impact on users, it is important that the impacts of these behaviors continue to be evaluated in different settings, as the nature and degree of the harm caused may vary in different scenarios \emph{eriksson2025can}. Current works that measure the domain-specific impacts of sycophancy have primarily focused on opinion-conforming behaviors, excess praising, and a lack of pushback. There are many sycophantic behaviors and application spaces that have yet to be studied, and the medium- to long-term impacts of sycophancy in general have been underexplored.
% \end{enumerate}

% \subsection{Use updated evaluations to inform development and deployment of targeted mitigation strategies}
% Several works [CITE] have proposed mitigation strategies that may reduce specific people-pleasing or sycophantic behaviors in LLMs. These benchmarks may work in isolated settings or for the specific sycophantic behaviors tested; however, with our current benchmarks, it is difficult to assess the downstream impacts these mitigations can have on real users and institutions. There may also be not-yet-explored or quantified sycophantic behaviors that current mitigation strategies may or may not transfer to. Steps 1 through 5 of our framework allow researchers to better assess these impacts and selectively deploy targeted mitigations for specific application scenarios where certain sycophantic behaviors have been linked to negative outcomes. 

% [TODO: examples - the verified negative impacts of social sycophancy behaviors should lead to more research, benchmarks, and targeted mitigation strategies surrounding these behaviors]

% [TODO: list some specific mitigation strategies that have been introduced in prior work. Mention that these have mostly been targeted towards opinion conformity behaviors, with a particular focus on improving accuracy, and that it is important to test the effects of these mitigation strategies on other sycophancy benchmarks and downstream impacts. Mitigation strategies should also be developed which target other sycophantic behaviors.]
% \section{Important Final Considerations}
% \subsection{Designing evaluations based on actual user behavior - should shift from focusing mainly on simulated, one-to-two-turn interactions to longer multi-turn interactions, and studies should cover the variety of ways in which users use LLMs}

% \subsection{Equity: sycophancy may harm some communities, minoritized groups, or populations underrepresented in the training data. For instance, homogenization of writing styles to a more Western style or younger users whose brains are still developing. This should be carefully studied}

% \subsection{Importance of studying downstream effects at both a macro and a micro level}
\section{Conclusion}

In this work, we proposed a unified construct of AI sycophancy, where sycophantic behaviors in AI systems are characterized by displacement of one or more context-relevant norms associated with observed approval-seeking behavior in model generations. This definitions allows researchers to empirically determine which behaviors should be considered sycophantic rather than attempt to redefine the concept, reconcile divergent definitions, or anthropomorphize AI systems. 



% \begin{enumerate}
%     \item Motivation: findings related to people-pleasing in LLMs and case studies showing the harms of sycophancy
%     \item Trends in AI research, news coverage, and online discourse
%     \item Position: how researchers measure AI sycophancy often diverges from users' perceptions of AI sycophancy and major concerns about it.
%     \item Position: We are at the point now where we have largely empirically validated the existence of sycophancy in LLMs, and even developed some mitigation strategies for it
%     \item Position: now, we need more research on the long-term impacts of sycophancy 
%     \begin{enumerate}
%         \item Can help us better understand the dangers of this phenomenon, and which sycophantic behaviors causes the most harm, so we can target our mitigation strategies
%         \item Can also help us distinguish sycophancy from things like personalization and empathetic behavior - the line between adjusting for human preferences and excessive people-pleasing could be drawn based on the negative impacts over time
%     \end{enumerate}
%     \item One way to distinguish sycophancy from normal preference optimization, in the case of a lack of ground truth, could be by studying the impacts of this behavior (characterize ``excessive'')
%     \item Personalization vs. sycophancy
%     \item Empathetic behavior vs. sycophancy
%     \item Describe gap in the literature and need for a survey paper that proposes an al-encompassing definition of AI sycophancy and reviews works that span different domains
%     \item Brief timeline of AI sycophancy research
%     \item Introduce review and describe methods used
%     \item Taxonomy of AI sycophancy research 
%     \begin{enumerate}
%         \item Different disciplines
%         \item Modalities: vision, text, audio
%         \item Empirical methods
%         \item Simulations vs. user studies
%         \item Single-turn vs. multi-turn
%     \end{enumerate}
%     \item Metrics used to measure sycophancy
%     \begin{enumerate}
%         \item Switching behavior
%         \item Divergence from ground truth
%         \item Excessive flattery
%         \item Factuality
%         \item Bayesian rationality
%     \end{enumerate}
%     \item Findings
%     \begin{enumerate}
%         \item Sycophantic behaviors identified
%          \item User perception of AI sycophancy and impact on users' judgments
%          \item Impact on multi-agent interactions
%         \begin{enumerate}
%             \item Controlled experiments
%             \item Surveys
%             \item Social media data
%         \end{enumerate}
%         \item Mitigation strategies
%     \end{enumerate}
%     \item Gaps in existing research
%     \item Lack of standardized definition
%     \item Future directions
%     \begin{enumerate}
%     \item More studies examining the effects of sycophantic behavior on users
%         \item More longitudinal studies
%         \item Multi-turn conversations and conversation length vs. sycophancy observed
%         \item Sycophancy in-the-wild
%     \end{enumerate}
% \end{enumerate}


% \section{Experiments/Analyses}

% \begin{enumerate}
%     \item Keyword searches for terms related to AI sycophancy in CS venues (ACM, IEEE, ACL, etc.) and in other disciplines (psychology, ethics, philosophy, etc.)
%     \begin{enumerate}
%         \item ProQuest
%         \item PhilPapers
%         \item Academic Search Ultimate
%         \item JSTOR
%         \item OneSearch - can create a user profile and save searches to user profile. 
%         Look at the subject breakdown and filter by language (English) and maybe filter out reviews. Can see the breakdown in disciplines
%         \item May have to do individual databases for abstract searches
%         \item Philosopher's index
%         \item Project Muse
%         \item ProQuest One Academic (dissertations and theses)
%         \item News/Media: Access World News Research Collection (Newsbank), New York Times, Washington Post, Wall Street Journal
%         \item Web of Science - good for citation tracking, data visualizations, etc. Also has abstract search
%         \item Look at OneSearch and see which disciplines and media types are coming up then target my search further
%     \end{enumerate}
%     \item Small Reddit experiment analyzing posts on r/ChatGPT and other subreddits about generative AI models

% \end{enumerate}

\bibliography{anthology-1,anthology-2,custom}

\appendix

\section{Example Appendix}
\label{sec:appendix}

This is an appendix.

\end{document}
